针对移动金融场景中高发的语义欺诈问题，传统手机银行风控存在语义解析能力不足、用户行为感知维度单一、AI决策缺乏可审计性等缺陷。本文依托AI Agent与MCP标准化交互协议，设计并实现了一套HIRD四层分层风控系统。系统由感知集成层（H）、认知决策层（I）、规则治理层（R）、业务执行层（D）组成，遵循“数据单向流转、权限逐级收敛”的设计准则，从架构层面隔离AI分析权限与业务执行权限，确保Agent“只析不执”。在感知层，基于认知负荷理论构建了包含有效输入时长、非预期停顿次数、内容修改频次、页面停留时长、操作节奏方差在内的五类微行为特征指标体系，通过Flutter框架实现移动端无感采集与标准化快照封装。在认知层，结合检索增强生成技术，由AI Agent完成交易文本的意图级语义解析与欺诈案例匹配。在规则治理层，设计了场景自适应动态权重算法，将大模型概率性输出收敛为确定性风控指令，并通过三级阈值路由实现放行、二次核验、拦截的分级处置与人工复核兜底。在工程实现上，系统基于FastAPI框架完成后端异步编排与双路径决策分流，基于MCP协议搭建标准化数据交互中枢，实现原子工具封装、幂等校验与全链路审计。测试结果表明，系统风控预检时延控制在100ms以内，语义欺诈识别召回率达94.2%，误报率控制在1.15%，敏感信息脱敏率达100%，可有效识别隐性社交工程欺诈，满足金融场景实时性、安全性与合规审计要求。